\documentclass[10pt,a4paper]{amsart}
\usepackage{amsmath,amssymb,amsthm,graphicx}
\usepackage{verbatim}
\newenvironment{reference}{\comment}{\endcomment}
%\newenvironment{reference}{}{}
\newenvironment{slogan}{\comment}{\endcomment}
\newenvironment{history}{\comment}{\endcomment}
% For commutative diagrams we use Xy-pic
\usepackage[all]{xy}

% We use 2cell for 2-commutative diagrams.
\xyoption{2cell}
% Theorem environments.
%
\theoremstyle{plain}
\newtheorem{theorem}[subsection]{Theorem}
\newtheorem{proposition}[subsection]{Proposition}
\newtheorem{lemma}[subsection]{Lemma}

\theoremstyle{definition}
\newtheorem{definition}[subsection]{Definition}
\newtheorem{example}[subsection]{Example}
\newtheorem{exercise}[subsection]{Exercise}
\newtheorem{situation}[subsection]{Situation}

\theoremstyle{remark}
\newtheorem{remark}[subsection]{Remark}
\newtheorem{remarks}[subsection]{Remarks}

\numberwithin{equation}{subsection}

% Macros
%
\def\lim{\mathop{\mathrm{lim}}\nolimits}
\def\colim{\mathop{\mathrm{colim}}\nolimits}
\def\Spec{\mathop{\mathrm{Spec}}}
\def\Hom{\mathop{\mathrm{Hom}}\nolimits}
\def\Ext{\mathop{\mathrm{Ext}}\nolimits}
\def\SheafHom{\mathop{\mathcal{H}\!\mathit{om}}\nolimits}
\def\SheafExt{\mathop{\mathcal{E}\!\mathit{xt}}\nolimits}
\def\Sch{\mathit{Sch}}
\def\Mor{\mathop{\mathrm{Mor}}\nolimits}
\def\Ob{\mathop{\mathrm{Ob}}\nolimits}
\def\Sh{\mathop{\mathit{Sh}}\nolimits}
\def\NL{\mathop{N\!L}\nolimits}
\def\CH{\mathop{\mathrm{CH}}\nolimits}
\def\proetale{{pro\text{-}\acute{e}tale}}
\def\etale{{\acute{e}tale}}
\def\QCoh{\mathit{QCoh}}
\def\Ker{\mathop{\mathrm{Ker}}}
\def\Im{\mathop{\mathrm{Im}}}
\def\Coker{\mathop{\mathrm{Coker}}}
\def\Coim{\mathop{\mathrm{Coim}}}

% Boxtimes
%
\DeclareMathSymbol{\boxtimes}{\mathbin}{AMSa}{"02}

%
% Macros for moduli stacks/spaces
%
\def\QCohstack{\mathcal{QC}\!\mathit{oh}}
\def\Cohstack{\mathcal{C}\!\mathit{oh}}
\def\Spacesstack{\mathcal{S}\!\mathit{paces}}
\def\Quotfunctor{\mathrm{Quot}}
\def\Hilbfunctor{\mathrm{Hilb}}
\def\Curvesstack{\mathcal{C}\!\mathit{urves}}
\def\Polarizedstack{\mathcal{P}\!\mathit{olarized}}
\def\Complexesstack{\mathcal{C}\!\mathit{omplexes}}
% \Pic is the operator that assigns to X its picard group, usage \Pic(X)
% \Picardstack_{X/B} denotes the Picard stack of X over B
% \Picardfunctor_{X/B} denotes the Picard functor of X over B
\def\Pic{\mathop{\mathrm{Pic}}\nolimits}
\def\Picardstack{\mathcal{P}\!\mathit{ic}}
\def\Picardfunctor{\mathrm{Pic}}
\def\Deformationcategory{\mathcal{D}\!\mathit{ef}}
\usepackage{cleveref}
\crefname{originallemma}{original lemma}{original lemmas}
\Crefname{originallemma}{Original lemma}{Original lemmas}
\makeatletter
\newlabel{lemma-base-change-quasi-compact}{{lemma-base-change-quasi-compact}{source}{}{sourcepointer-1}{}}
\newlabel{lemma-base-change-quasi-compact@cref}{{[originallemma][0][]lemma-base-change-quasi-compact}{[1][0][]source}{}{}{}}
\makeatother
\crefname{originallemma}{original lemma}{original lemmas}
\Crefname{originallemma}{Original lemma}{Original lemmas}
\makeatletter
\newlabel{lemma-base-change-separated}{{lemma-base-change-separated}{source}{}{sourcepointer-2}{}}
\newlabel{lemma-base-change-separated@cref}{{[originallemma][0][]lemma-base-change-separated}{[1][0][]source}{}{}{}}
\makeatother
\crefname{originallemma}{original lemma}{original lemmas}
\Crefname{originallemma}{Original lemma}{Original lemmas}
\makeatletter
\newlabel{lemma-compose-after-separated}{{lemma-compose-after-separated}{source}{}{sourcepointer-3}{}}
\newlabel{lemma-compose-after-separated@cref}{{[originallemma][0][]lemma-compose-after-separated}{[1][0][]source}{}{}{}}
\makeatother
\crefname{originallemma}{original lemma}{original lemmas}
\Crefname{originallemma}{Original lemma}{Original lemmas}
\makeatletter
\newlabel{lemma-composition-quasi-compact}{{lemma-composition-quasi-compact}{source}{}{sourcepointer-4}{}}
\newlabel{lemma-composition-quasi-compact@cref}{{[originallemma][0][]lemma-composition-quasi-compact}{[1][0][]source}{}{}{}}
\makeatother
\crefname{originallemma}{original lemma}{original lemmas}
\Crefname{originallemma}{Original lemma}{Original lemmas}
\makeatletter
\newlabel{lemma-quasi-compact-permanence}{{lemma-quasi-compact-permanence}{source}{}{sourcepointer-5}{}}
\newlabel{lemma-quasi-compact-permanence@cref}{{[originallemma][0][]lemma-quasi-compact-permanence}{[1][0][]source}{}{}{}}
\makeatother
\crefname{originallemma}{original lemma}{original lemmas}
\Crefname{originallemma}{Original lemma}{Original lemmas}
\makeatletter
\newlabel{lemma-separated-over-separated}{{lemma-separated-over-separated}{source}{}{sourcepointer-6}{}}
\newlabel{lemma-separated-over-separated@cref}{{[originallemma][0][]lemma-separated-over-separated}{[1][0][]source}{}{}{}}
\makeatother
\title{Quasi-compactness of the source forces quasi-compactness of the morphism, and fibre products of quasi-compact quasi-separated stacks stay so}
\author{The Stacks Project Authors}
\date{stacks-project:stacks-morphisms.tex:local-snapshot-20261012; GNU FDL 1.2 or later}
\begin{document}
\maketitle
\noindent\textit{Editorial note.} All delivered mathematical text is the author's own, copied from the pinned Stacks Project chapter source sources/stacks-algebraic-stacks-morphisms/pinned-original-source.tex; no statement or proof step was rewritten, reordered or removed.\par\medskip
\noindent\textit{Editorial note.} Preamble provenance: the Stacks Project keeps one single shared preamble for every chapter; this card takes it directly from the staged sources/stacks-algebraic-stacks-morphisms/preamble.tex as three verbatim spans (lines 8-12, 14-18 and 154-217), so no derived file was written and nothing was added or redefined. The wrapper class and the base packages are supplied by the assembler.\par\medskip
\noindent\textit{Editorial note.} Bibliography provenance: neither the counted statement nor its proof contains a citation (this lane rescans the whole span before writing the card), so the card carries no bibliography and no entry was reconstructed.\par\medskip
\noindent\textit{Editorial note.} Selection provenance: this card was opened under the file-based control contract (CONTROL.md), with the mechanical pre-card gate run on its statement and proof interval before the card was written.\par\medskip
\noindent\textit{Allowed standard prerequisites (editorial).}\begin{itemize}
\item The statement that a morphism with quasi-compact source and quasi-separated target is quasi-compact. Invoked at the start of the counted proof (source line 2081). Proves (1).
\item The statement that composing a separated morphism after a morphism preserves quasi-separatedness. Invoked inside the counted proof (source lines 2082-2084). Upgrades (1) to (2).
\item The base change permanence of quasi-compactness and of the separation properties. Invoked inside the counted proof (source lines 2088-2092). Makes the projection of the fibre product quasi-compact and quasi-separated.
\item The statements comparing a stack with its structure morphism over a quasi-compact quasi-separated base. Invoked at the end of the counted proof (source lines 2093-2097). Concludes that the fibre product itself is quasi-compact and quasi-separated.
\end{itemize}\medskip
\ifdefined\hypertarget\hypertarget{sourcepointer-1}{}\fi
\noindent\textit{Source reference.} Original lemma lemma-base-change-quasi-compact: sources/stacks-algebraic-stacks-morphisms/pinned-original-source.tex, source lines 1984--1984. Navigation pointer only: the counted text refers to the Stacks item lemma-base-change-quasi-compact, which is an existing labeled result of this same chapter and is not part of this card. The pointer carries the target's own label as its number and its definition line as its span; its content is not reproduced here.\par\medskip
\ifdefined\hypertarget\hypertarget{sourcepointer-2}{}\fi
\noindent\textit{Source reference.} Original lemma lemma-base-change-separated: sources/stacks-algebraic-stacks-morphisms/pinned-original-source.tex, source lines 703--703. Navigation pointer only: the counted text refers to the Stacks item lemma-base-change-separated, which is an existing labeled result of this same chapter and is not part of this card. The pointer carries the target's own label as its number and its definition line as its span; its content is not reproduced here.\par\medskip
\ifdefined\hypertarget\hypertarget{sourcepointer-3}{}\fi
\noindent\textit{Source reference.} Original lemma lemma-compose-after-separated: sources/stacks-algebraic-stacks-morphisms/pinned-original-source.tex, source lines 929--929. Navigation pointer only: the counted text refers to the Stacks item lemma-compose-after-separated, which is an existing labeled result of this same chapter and is not part of this card. The pointer carries the target's own label as its number and its definition line as its span; its content is not reproduced here.\par\medskip
\ifdefined\hypertarget\hypertarget{sourcepointer-4}{}\fi
\noindent\textit{Source reference.} Original lemma lemma-composition-quasi-compact: sources/stacks-algebraic-stacks-morphisms/pinned-original-source.tex, source lines 1994--1994. Navigation pointer only: the counted text refers to the Stacks item lemma-composition-quasi-compact, which is an existing labeled result of this same chapter and is not part of this card. The pointer carries the target's own label as its number and its definition line as its span; its content is not reproduced here.\par\medskip
\ifdefined\hypertarget\hypertarget{sourcepointer-5}{}\fi
\noindent\textit{Source reference.} Original lemma lemma-quasi-compact-permanence: sources/stacks-algebraic-stacks-morphisms/pinned-original-source.tex, source lines 2044--2044. Navigation pointer only: the counted text refers to the Stacks item lemma-quasi-compact-permanence, which is an existing labeled result of this same chapter and is not part of this card. The pointer carries the target's own label as its number and its definition line as its span; its content is not reproduced here.\par\medskip
\ifdefined\hypertarget\hypertarget{sourcepointer-6}{}\fi
\noindent\textit{Source reference.} Original lemma lemma-separated-over-separated: sources/stacks-algebraic-stacks-morphisms/pinned-original-source.tex, source lines 888--888. Navigation pointer only: the counted text refers to the Stacks item lemma-separated-over-separated, which is an existing labeled result of this same chapter and is not part of this card. The pointer carries the target's own label as its number and its definition line as its span; its content is not reproduced here.\par\medskip
\begin{lemma}
\label{lemma-quasi-compact-quasi-separated-permanence}
Let $f : \mathcal{X} \to \mathcal{Y}$ be a morphism of algebraic stacks.
\begin{enumerate}
\item If $\mathcal{X}$ is quasi-compact and $\mathcal{Y}$ is
quasi-separated, then $f$ is quasi-compact.
\item If $\mathcal{X}$ is quasi-compact and quasi-separated and $\mathcal{Y}$
is quasi-separated, then $f$ is quasi-compact and quasi-separated.
\item A fibre product of quasi-compact and quasi-separated algebraic stacks
is quasi-compact and quasi-separated.
\end{enumerate}
\end{lemma}

\begin{proof}
Part (1) follows from
Lemma \ref{lemma-quasi-compact-permanence}.
Part (2) follows from (1) and
Lemma \ref{lemma-compose-after-separated}.
For (3) let $\mathcal{X} \to \mathcal{Y}$ and $\mathcal{Z} \to \mathcal{Y}$
be morphisms of quasi-compact and quasi-separated algebraic stacks.
Then $\mathcal{X} \times_\mathcal{Y} \mathcal{Z} \to \mathcal{Z}$
is quasi-compact and quasi-separated as a base change of
$\mathcal{X} \to \mathcal{Y}$ using (2) and
Lemmas \ref{lemma-base-change-quasi-compact} and
\ref{lemma-base-change-separated}.
Hence $\mathcal{X} \times_\mathcal{Y} \mathcal{Z}$
is quasi-compact and quasi-separated as
an algebraic stack quasi-compact and quasi-separated over
$\mathcal{Z}$, see
Lemmas \ref{lemma-separated-over-separated} and
\ref{lemma-composition-quasi-compact}.
\end{proof}

\end{document}
